\providecommand{\method}{\ours}

\section{Methodology}
\label{sec:method}

\method is designed to be practical, accurate, and efficient. It has to run
with the models a team already uses and leave stored records intact, it has to
keep a name, a date, a count, or the distinction between two events recoverable
after writing and after consolidation, and it has to answer from a bounded
evidence set however long the history grows. This section describes the
layered design that delivers the three.

\method separates storage from indexing (\Cref{fig:mirix-overview}). A \emph{typed evidence plane} stores complete
memories and is the system of record (\Cref{sec:typed-evidence}). A
\emph{provenance hypergraph} indexes the entities, events, relations, roles, and
times in those memories, and every fact in it cites the records it came from
(\Cref{sec:hypergraph}). Three paths operate on the two layers, and
\Cref{alg:mirix} states each of them in pseudocode. The write path stores a
memory and then indexes it. The read path plans a query, traverses the
graph, and dereferences only the records it selects (\Cref{sec:retrieval}).
AutoDream, the third path, runs outside the other two (\Cref{sec:autodream}). On semantic memory it
merges duplicate anchors in the hypergraph and returns the evidence store
unchanged, so semantic consolidation changes how evidence is connected, never
what evidence exists. On procedural
memory, AutoDream instead distills task sessions into skills that cite the
experiences behind them. \Cref{app:method-invariants} collects the properties that hold regardless of
how the language models behave.

\subsection{Problem setup}
\label{sec:method-setup}

For a user \(u\), the system observes a stream
\(\mathcal X_u=(x_1,\ldots,x_T)\) of conversational or multimodal inputs. When it
writes \(x_t\), future queries are unknown, so the writer must produce a
persistent state
\begin{equation}
    \mathcal S_u^t=(P_u^t,G_u^t)
\end{equation}
from which a later query \(q\) can be answered with a bounded evidence set
\(\mathcal R(q;\mathcal S_u^t)\). The evidence store \(P_u\) collects the typed
stores, among them the episodic store \(P_u^E\), the semantic store \(P_u^S\),
and the procedural store \(P_u^P\). We seek high answer utility under three
technical requirements.
\begin{enumerate}[leftmargin=*,nosep]
    \item \textbf{Source fidelity.} Concrete evidence remains recoverable after
    writing and after consolidation.
    \item \textbf{Bounded access.} A query inspects a small graph neighborhood
    and a bounded set of cited rows instead of scanning the full history.
    \item \textbf{Temporal identity.} Repeated descriptions of one event may
    share evidence, but distinct occurrences must not collapse merely because
    they are worded alike.
\end{enumerate}
Source fidelity and temporal identity are what make the system accurate and
its answers traceable, and bounded access is what keeps it efficient. The rest
of this section shows how the two layers meet these requirements.
Each subsection explains a block of \Cref{alg:mirix}, and the numbers on its right
point to the equations that define each step.

\begin{algorithm}[t]
\caption{\method for one user \(u\). Calls marked \(\star\) use a language model.}
\label{alg:mirix}
\small
\algrenewcommand{\algorithmiccomment}[1]{\hfill#1}
\begin{minipage}[t]{0.485\textwidth}
\begin{algorithmic}[1]
\Procedure{Write}{$x_t$}
  \For{\(k\in R_\theta^\star(x_t)\)}\label{alg:route}
    \State \(m\gets W_k^\star(x_t,P_u)\), \(\textsc{Commit}(P_u,m)\)\Comment{\eqref{eq:memory-record}}\label{alg:commit}
    \If{\(k\in\{\mathrm{episodic},\mathrm{semantic}\}\)}
      \State \(\Phi\gets\textsc{Extract}^\star(m)\cup\textsc{Repair}^\star(m)\)\Comment{\eqref{eq:frame}}\label{alg:extract}
      \State \(A_m\gets\textsc{Admit}(\textsc{Canon}^\star(\Phi))\)
      \State \(I_a\gets I_a\cup\{i_m\}\) for \(a\in A_m\)
      \For{\(f\in\Phi\) with \(|A_f|\ge 2\)}
        \State \(\textsc{Upsert}(G_u,\operatorname{id}(f),f)\)\Comment{\eqref{eq:semantic-frame-id}, \eqref{eq:episodic-frame-id}}\label{alg:upsert}
        \State \(C_f\gets C_f\cup\{i_m\}\)
      \EndFor
    \EndIf
  \EndFor
\EndProcedure
\Procedure{Read}{$q$}
  \State \(\Pi\gets\textsc{Plan}(q)\)\Comment{\eqref{eq:query-plan}}\label{alg:plan}
  \State \(A_q\gets\textsc{Top}(\textsc{Recall}(q),\,s_A,\,K_A)\)\Comment{\eqref{eq:anchor-score}}
  \State \(F_q\gets\textsc{Top}(\textsc{Adj}(A_q,\mathcal P_q),\,s_F,\,B_q^R)\)\Comment{\eqref{eq:fact-score}}
  \State \Return \(F_q\), \(\textsc{Fuse}(\textsc{Rows}(A_q),\,\bigcup_{f\in F_q}C_f;\,B_q^B,B_q^R)\)
\EndProcedure
\algstore{mirixalg}
\end{algorithmic}
\end{minipage}\hfill
\begin{minipage}[t]{0.485\textwidth}
\begin{algorithmic}[1]
\algrestore{mirixalg}
\Procedure{Answer}{$q$}
  \State \(\hat y\gets\textsc{Agent}^\star(q)\) with at most five \textsc{Read} calls
  \If{\(\textsc{Check}(\hat y)\)}
    \State \(\hat y\gets\textsc{Rewrite}^\star(\hat y,\textsc{Read}(q'))\)\label{alg:retry}
  \EndIf
  \State \Return \(\hat y\)
\EndProcedure
\Procedure{SemanticDream}{$\Delta_u$}\Comment{\eqref{eq:dream-contract}}
  \For{\((a,b)\in\mathcal C_\Delta\setminus\mathcal N\)}\Comment{\eqref{eq:dream-candidates}}\label{alg:dream-pairs}
    \If{\(\textsc{Same}^\star(a,b)\)}
      \State \(\textsc{Merge}(a\to b)\)
    \Else
      \State \(\mathcal N\gets\mathcal N\cup\{(a,b)\}\)
    \EndIf
  \EndFor
\EndProcedure
\Procedure{ProceduralDream}{$S_u$}
  \State \(e_{1:m}\gets\textsc{Distill}^\star(S_u)\)
  \State \(P_u^P\gets\textsc{Curate}^\star(P_u^P,e_{1:m})\)\label{alg:procedural}
\EndProcedure
\end{algorithmic}
\end{minipage}
\end{algorithm}

\subsection{Typed evidence plane}
\label{sec:typed-evidence}

\paragraph{Multi-agent routing.}
A Meta Memory Manager routes each input \(x_t\) to the subset
\(\mathcal K_t=R_\theta(x_t)\) of memory types it updates (\Cref{alg:mirix},
line~\ref{alg:route}). Each of the six types, Core, Episodic, Semantic,
Procedural, Resource, and Knowledge Vault, has its own writer, and the writer
of a selected type \(k\) produces the record \(m_t^{(k)}=W_k(x_t,P_u^{t-1})\)
and owns the schema and update operations of that type. Only episodic and
semantic records are indexed by the hypergraph. \Cref{app:method-routing} gives
the roles of the six types and the reason for this split.

\paragraph{Evidence records.}
Each of the six memory types keeps its records in its own schema. The hypergraph
indexes the episodic and semantic records, and each of them is represented as
\begin{equation}
    m=(i,k,y,\tau_e,\tau_m,\rho,\mathbf z),
    \label{eq:memory-record}
\end{equation}
where \(i\) is an immutable row identifier,
\(k\in\{\mathrm{episodic},\mathrm{semantic}\}\) is its memory type, \(y\) is the
stored text, \(\tau_e\) is the event time when one applies, \(\tau_m\) is the
mention or write time, \(\rho\) records source and chunk provenance, and
\(\mathbf z\) holds the search embeddings. Full summaries and details remain in
the vector database. The graph stores only identifiers and compact structural metadata.

Writers keep concrete details such as names, counts, and dates in each record,
and a record is committed to the vector database before graph construction
begins (line~\ref{alg:commit}), so a failed graph write never rolls back a
stored memory (\Cref{app:method-routing}).

\subsection{A provenance-preserving n-ary hypergraph}
\label{sec:hypergraph}

The second half of the write path turns each episodic or semantic record into
graph structure (\Cref{alg:mirix}, lines~\ref{alg:extract}--\ref{alg:upsert}).
For each user, \method maintains a hypergraph
\begin{equation}
    G_u=(\mathcal A_u,\mathcal F_u,\mathcal E_u)
\end{equation}
with entity or concept anchors \(\mathcal A_u\), reified fact frames
\(\mathcal F_u\), and role-typed incidence edges \(\mathcal E_u\).
\Cref{app:method-implementation} shows the hypergraph \method builds for one LoCoMo conversation.

\paragraph{Anchors.}
An anchor
\begin{equation}
    a=(i_a,n_a,t_a,\mathbf e_a,I_a^E,I_a^S)
\end{equation}
stores a canonical name \(n_a\), an entity type \(t_a\), a name embedding
\(\mathbf e_a\), and the identifiers \(I_a^E\) and \(I_a^S\) of the episodic and
semantic rows that mention it. A deterministic key \(\kappa(n)\) normalizes
case, punctuation, function words, word order, and singular and plural variants
of a name while keeping numbers, and a deterministic admission gate rejects
generic filing terms and dialogue-role pseudo-entities and caps the anchors
admitted from one record (\Cref{app:method-graph}).

\paragraph{Registry-guided canonicalization.}
Surface normalization cannot decide whether two semantically close names denote
the same entity. For every extracted entity, approximate nearest-neighbor search
therefore retrieves a few existing anchors of the user above a similarity
threshold, and one language-model call per record makes a conservative
same-entity decision for all of its entities. A rename is accepted only when it
selects a retrieved candidate, so embeddings propose comparisons but never
authorize a merge on their own. The neighbor count and the threshold, with their relation to the vector-index score, are in \Cref{app:method-graph}.

\paragraph{N-ary frames.}
For each episodic or semantic record, an extractor produces a set of frames
\begin{equation}
    f=
    \left(
        p,\,
        \{(r_j,a_j)\}_{j=1}^{d_f},\,
        \ell_f,\,
        \tau_e,\tau_m,\,
        C_f
    \right),
    \label{eq:frame}
\end{equation}
where \(p\) is a canonical predicate, \(r_j\) is the open-vocabulary role of
anchor \(a_j\), \(\ell_f\) maps literal attributes to their values, and \(C_f\)
is the set of identifiers of the vector-database records the frame cites. A flight, for example,
becomes one frame with agent, origin, destination, carrier, and companion arms
instead of several disconnected triples, and owns(Mia, Suica card) in
\Cref{fig:mirix-user-centric} is a frame that cites the May 2 record. The extractor is instructed to keep dates, prices, counts, durations, and
measurements as literals on the frame rather than as anchors, so that a value
stays attached to the assertion it qualifies. A frame is materialized only if
its arguments resolve to at least two distinct anchors of the user. When a deterministic audit finds that
the extraction misses high-signal source details, one bounded repair call may
add frames that the source supports, and it cannot remove anything the primary
extraction produced. \Cref{app:method-graph} describes the audit.

\paragraph{Occurrence identity.}
Semantic claims and episodic occurrences follow different identity rules. Let
\(\sigma(f)\) be the sorted set of role and anchor keys
\(\{r_j:\kappa(n_{a_j})\}\). A semantic frame is identified by
\begin{equation}
    \operatorname{id}_S(f)
    =H(u,p,\sigma(f)),
    \label{eq:semantic-frame-id}
\end{equation}
so repeated support for the same durable claim accumulates citations on one
node. An episodic frame adds an occurrence key,
\begin{align}
    \operatorname{id}_E(f)
        &=H(u,p,\sigma(f),\chi(f)), \label{eq:episodic-frame-id}\\
    \chi(f)
        &=
        \begin{cases}
        \text{explicit event time}, & \text{if extracted},\\
        \text{immutable memory id}, & \text{otherwise}.
        \end{cases}
\end{align}
Two mentions of the same dated event therefore share one frame, while two undated
events with similar wording, such as two runs on different days, stay separate. A
frame keeps the time of the event apart from the time it was mentioned, and whether
the event is planned or done does not change its identity.

\subsection{Graph-owned evidence retrieval}
\label{sec:retrieval}

\method treats retrieval as the construction of a bounded evidence set rather
than as open-ended similarity search. Episodic and semantic candidates must
enter through the graph, and the vector database dereferences and reranks only the
identifiers that graph traversal selects. This boundary keeps a hidden
full-table search from doing the work attributed to the graph. Only a graph
read that raises an error falls back to direct evidence retrieval. A query passes
through planning, anchor entry, expansion along two evidence lanes, and fusion,
and a verifier then checks the answer
(\Cref{alg:mirix}, lines~\ref{alg:plan}--\ref{alg:retry}). For the question in
\Cref{fig:mirix-user-centric}, the planner chooses the single operator, enters
the graph at Mia and Tokyo transit, and returns three facts with five rows.

\paragraph{Query plan.}
A deterministic planner maps a query to
\begin{equation}
    \Pi(q)=
    \left(
        o_q,\mathcal P_q,\mathcal C_q,\mathcal Y_q,\mathcal T_q,
        B_q^B,B_q^R
    \right),
    \label{eq:query-plan}
\end{equation}
where the operator \(o_q\) is single, temporal, intersection, list-union, or
count, \(\mathcal P_q\) holds canonical relation families and their surface
forms, \(\mathcal C_q\) holds non-relational constraints, \(\mathcal Y_q\) holds
expected answer types, and \(\mathcal T_q\) holds date and temporal-order
constraints. The two budgets cap the base-recall rows and the relation-cited
rows, and \(B_q^R\) also caps the graph facts that the relation lane renders.
The budgets grow with the operator, from 6 base rows and 4 cited rows for a
single-value lookup to 8 and 12 for a count, since counting and set
construction need broader evidence.
\Cref{tab:lane-budgets} in \Cref{app:method-retrieval} lists all five
operators.

\paragraph{Anchor entry.}
The query embedding first recalls fifteen anchors by name. A relation query adds
the anchors that the question names exactly, and the union is scored by
\begin{equation}
    s_A(a,q)=s_{\mathrm{vec}}+s_{\mathrm{exact}}+
    s_{\mathrm{lex}}+s_{\mathrm{pred}}+s_{\mathrm{spec}}
    -s_{\mathrm{hub}}.
    \label{eq:anchor-score}
\end{equation}
The vector score dominates. Exact mention, token overlap, predicate
compatibility, and multi-token specificity break close ties, and a logarithmic
degree penalty suppresses generic hubs. The top fifteen anchors are kept, which
is the cap \(K_A\) of \Cref{alg:mirix}. The cap on exactly named anchors and the hub and predicate rules are in \Cref{app:method-retrieval}.

\paragraph{Two evidence lanes.}
The selected anchors open two bounded lanes. In the base-recall lane, anchors
yield the episodic and semantic identifiers directly associated with them, and
traversing an incident n-ary frame adds the frame's cited memories and a bounded
set of memories associated with its co-participants. In the relation-citation
lane, which opens when the query plan contains a relation, the exactly named anchors seed an indexed adjacency lookup over their incident
frames. Predicate filtering applies before a cap of 240 frames, the fact score
below ranks the rest, and at most \(B_q^R\) of them are rendered as explicit
graph facts that contribute their exact citations into the vector database. The provenance
recomputation of the base lane and the object filter of the relation lane are
in \Cref{app:method-retrieval}.

\paragraph{Fact scoring.}
Each relation-lane frame receives the additive score
\begin{align}
    s_F(f,q)={}&
    s_{\mathrm{rel}}+s_{\mathrm{entity}}+s_{\mathrm{constraint}}
    +s_{\mathrm{type}} \notag\\
    &+s_{\mathrm{time}}+s_{\mathrm{intersection}}
    +s_{\mathrm{specificity}} \notag\\
    &+s_{\mathrm{role}}+s_{\mathrm{state}}.
    \label{eq:fact-score}
\end{align}
The first seven terms score relation, named-entity, literal, answer-type, and
temporal compatibility, a shared object for intersection questions, and the
number of arguments a frame binds. The role term \(s_{\mathrm{role}}\) rewards agreement
between the argument direction implied by the question and the frame's typed
roles, and it penalizes incompatible roles. The state term \(s_{\mathrm{state}}\)
rewards an event time that matches a requested date and penalizes planned
occurrences in count, list-union, and temporal questions.
\Cref{app:method-retrieval} describes its fallback and the treatment of intersection questions.

\paragraph{Citation-balanced fusion.}
Rows in each lane are reranked, and fusion takes up to \(B_q^R\) cited rows,
then up to \(B_q^B\) base rows, and fills the remaining capacity from the rest
of either lane without duplicating row identifiers. The returned evidence
therefore contains both the graph facts and the source rows that justify them.
The reranking prior, the per-type split, and the rule for near-duplicate paraphrases are also in \Cref{app:method-retrieval}.

\paragraph{Answer-time verification.}
In our harness, the answering agent may call the retriever up to five times.
After it drafts an answer, a rule-based check requests one more search and one
rewrite by the answering model when the draft shows explicit uncertainty, lacks a value the question asks for,
or disagrees with the evidence ledger, for example a count without a number or a
month that neither the evidence nor the question contains. It also fires on
every intersection question (line~\ref{alg:retry}). The check runs in the answering agent rather than in
the memory server, and it never inspects a reference answer.
\Cref{app:method-verification} lists every trigger and the query of the extra
search.

\subsection{AutoDream: consolidation outside the write path}
\label{sec:autodream}

AutoDream runs outside the critical write path
(\Cref{alg:mirix}, lines~\ref{alg:dream-pairs}--\ref{alg:procedural}). It treats
the two kinds of memory it consolidates in opposite ways. On semantic memory it
is an explicit service call that never schedules itself, and it refines
identities under the contract
\begin{equation}
    \operatorname{Dream}(P_u,G_u;\Delta_u)
    = (P_u,\widetilde G_u),
    \label{eq:dream-contract}
\end{equation}
where \(\Delta_u\) is the newly written semantic frontier and \(\widetilde G_u\)
is the consolidated hypergraph. The evidence store \(P_u\) is returned
unchanged. On procedural memory it starts in the background once five task
sessions are pending, and it turns them into skills.

\paragraph{Semantic frontier.}
Each call receives the exact source-chunk identifiers of the batch just
written, which resolve to the newly written semantic rows. These rows define a
small frontier \(\Delta_u\), so the cost of consolidation depends on new
activity rather than on the total size of memory. The invocation schedule of our protocol is in \Cref{app:method-autodream}.

\paragraph{Candidate generation.}
Only anchors that are backed by semantic rows and belong to an identity-stable
type \(t\in\mathcal T_{\mathrm{safe}}=\{\text{person, organization, location,
object, concept}\}\) are eligible. For each eligible anchor \(a\in\mathcal A_\Delta\) backed by a frontier row, AutoDream takes the four nearest entries of the name index, which include \(a\) itself,
and forms the candidate pairs
\begin{equation}
    \mathcal C_\Delta=
    \left\{
    (a,b):
    \begin{array}{l}
    a\in\mathcal A_\Delta,\;
    b\in\operatorname{NN}_4(a)\setminus\{a\},\\
    \cos(\mathbf e_a,\mathbf e_b)\ge 0.86,\\
    t_a=t_b\in\mathcal T_{\mathrm{safe}}
    \end{array}
    \right\}.
    \label{eq:dream-candidates}
\end{equation}
Episodic-only anchors are excluded by construction. An LLM decides whether each
pair denotes the same entity and, if it does, chooses one of the two existing
names as the canonical representative. Its output is constrained to the
submitted pair, and invented representatives are rejected. In
\Cref{fig:mirix-agent-centric}, \emph{steel bottle} and \emph{stainless steel
bottle} form such a pair.

\paragraph{Bounded schedule.}
Every Dream round verifies at most a fixed number of unseen pairs, and the
final round also considers the eligible anchors that no earlier round used as a
source. Explicit negative decisions are cached as the set \(\mathcal N\) of
\Cref{alg:mirix}, so later cycles do not pay again to reconsider the same pair.
\Cref{app:method-autodream} states the caps of the intermediate, final, and
repair rounds and the cache key.

\paragraph{Lossless structural merge.}
For a confirmed merge \(a\rightarrow b\), the system unions the episodic and
semantic row identifiers, keeps aliases, sums mention counts, keeps the maximum
admission score, and rewires every role-typed frame arm to \(b\). Union, sum, and maximum are commutative and
associative, so the order of repeated merges does not change the reachable
evidence. The dropped anchor is removed only after its values and edges have
been transferred.

Semantic Dream never rewrites the memories in the vector database, and it consolidates neither
episodic-only anchors nor fact nodes. The asymmetry is deliberate. Semantic
identity can be refined incrementally, whereas event occurrences and their
citations are kept as observations. AutoDream therefore makes the graph less
fragmented without creating a textual summary whose embedding could hide the
details of its inputs.

\paragraph{Procedural refinement.}
For procedural memory, AutoDream refines instead of merging. A distiller judges
the outcome of each pending session from quotes in its transcript and distills
experiences \(e_1,\ldots,e_m\), each marked as a lesson worth reusing or a
mistake worth avoiding. A curator~\citep{zhang2026ace} turns these experiences
into skills in the procedural store \(P_u^P\), and it can create, edit, or
delete a skill. An edit can cite the experiences behind it with a quote of the
text it adds, and a citation is recorded only when that quote checks out, so a
cited step of a skill can be traced through its experiences to the quoted
sessions they came from. In \Cref{fig:mirix-agent-centric}, one failed purchase
yields an experience to avoid and one to reuse, and the curator adds two steps
to the shopping routine that cite the first and keeps a step that the second
supports. The distiller's unclear verdict and the curator's non-edit verdicts are in \Cref{app:method-autodream}.
